
Our investigation tests the gradient competition hypothesis through a two-stage experimental design: first establishing that spurious feature dominance exists (H-E1), then testing whether this dominance arises from stronger gradient signals (H-M1).

\subsubsection{3.1 Overview}

The gradient competition hypothesis posits that spurious features dominate because they produce stronger gradient signals during training. If true, we should observe:

\begin{enumerate}
\item \textbf{Existence:} GradCAM attribution on spurious regions exceeds attribution on core regions early in training
\item \textbf{Mechanism:} Gradient norms from spurious-aligned samples exceed those from minority samples, especially in early epochs
\end{enumerate}

\subsubsection{3.2 Experiment 1: Spurious Feature Attribution (H-E1)}

\textbf{Objective:} Measure whether spurious features dominate core features early in training.

We use GradCAM to compute class-discriminative attributions at layer4 (final convolutional block) of ResNet-50. For each validation sample, we obtain a spatial attribution map indicating which image regions contribute to the classification decision.

We define spurious and core regions based on Waterbirds image structure using a spatial heuristic:
\begin{itemize}
\item \textbf{Spurious region:} Upper 60\% of image (background habitat)
\item \textbf{Core region:} Lower 40\% of image (bird location)
\end{itemize}

For each epoch t, we compute the attribution ratio R(t) as the sum of spurious attributions divided by the sum of core attributions across validation samples.

\textbf{Success criterion:} R(t) > 1.0 for t < 10 (spurious dominance before epoch 10).

\subsubsection{3.3 Experiment 2: Gradient Norm Analysis (H-M1)}

\textbf{Objective:} Test whether spurious features receive stronger gradient signals.

We partition training samples into two groups based on spurious correlation alignment:
\begin{itemize}
\item \textbf{Spurious-aligned:} Samples where label matches background type (e.g., waterbird on water)
\item \textbf{Minority:} Samples where label mismatches background (e.g., waterbird on land)
\end{itemize}

For each training epoch, we compute L2 gradient norms at layer4 for each group and calculate the gradient norm ratio $\rho$(t) = G\_spurious(t) / G\_minority(t).

\textbf{Hypothesis:} If gradient competition drives spurious dominance, then $\rho$(t) > 1.5 in early epochs (1-10).

\textbf{Falsification criterion:} If $\rho$(t) < 1.0, the mechanism is inverted---minority samples produce stronger gradients.

\subsubsection{3.4 Training Configuration}

Both experiments use Waterbirds v1.0 dataset, ResNet-50 (ImageNet pretrained), SGD optimizer (momentum=0.9, weight\_decay=1e-4), batch size 128, and 50 epochs. H-E1 uses learning rate 0.01; H-M1 uses learning rate 0.001 for stable gradient measurements. H-M1 runs across 3 seeds (42, 123, 456).

\subsubsection{3.5 Gate Structure}

\begin{table}[htbp]
\centering
\resizebox{\columnwidth}{!}{%
\begin{tabular}{llll}
\toprule
Hypothesis & Gate Type & Condition & Action if Fail \\
\midrule
H-E1 & MUST\_WORK & R(t) > 1.0 for t < 10 & Abort (no spurious dominance) \\
H-M1 & MUST\_WORK & $\rho$(t) > 1.5 for t $\in$ [1,10] & Pivot (mechanism incorrect) \\
\bottomrule
\end{tabular}
}
\end{table}
